\section{Related work}
\textbf{Concept bottlenecks.} CBMs~\cite{koh2020concept} introduced the
image$\rightarrow$concepts$\rightarrow$class factorization and the associated
concept-level interpretability and intervention. Our concept layer follows their
supervised, class-level attribute setup. Subsequent analyses, however, show that
concepts predicted from the whole image are not necessarily learned as intended
and can encode (``leak'') class information beyond the evidence they
name~\cite{margeloiu2021concept}, which motivates predicting each concept from a
localized region.

\textbf{Hierarchy in the concept space.}
\citet{pittino2023hierarchical} aggregate concepts hierarchically for explainable
fine classification, motivating grouping structure over concepts;
\citet{sun2024boosting} add supervised, hierarchical concept learning with
structural constraints between coarse and fine concepts. Hierarchical
prototypes~\cite{hase2019interpretable} pursue interpretable recognition through
an explicit top-down structure. These works inform our part-grounded and
coarse-to-fine concept design.

\textbf{Hierarchy in the label space and better mistakes.}
\citet{bertinetto2020making} formalize mistake severity through the LCA of the
predicted and true class in a semantic hierarchy and propose hierarchy-aware
training losses. \citet{karthik2021nocost} show that severity can be reduced
purely at test time, without retraining, by choosing the class of minimum
expected risk (Conditional Risk Minimization, CRM)---the mechanism we adopt for
our label-space axis.

\textbf{Positioning.} The survey of \citet{poeta2025concept} frames hierarchical
concept structure as an open direction in concept-based explainability, which our
two-axis study addresses directly.
